\phant{THẢO LUẬN}\label{sec:thaoluan}

\subsection{Hai lớp kiến thức của mã song song}

Kết quả cho thấy năng lực của LLM trong mã HPC phụ thuộc vào hai lớp kiến thức
khác nhau. Lớp thứ nhất là \emph{cú pháp và thư viện chuẩn}: viết đúng vòng lặp,
đúng chữ ký hàm, đúng cách dùng \code{std::async}, \code{std::execution::par}
hay \code{\#pragma omp parallel for}. Lớp thứ hai là \emph{kiến thức hệ sinh
thái}: trình biên dịch nào, backend nào, phiên bản thư viện nào, tham số nào thực
sự có hiệu lực, và chi phí nào sẽ xuất hiện khi chạy trên phần cứng thật.

Tỉ lệ biên dịch cao (100\% ở ba mô hình) cho thấy lớp thứ nhất đã khá tốt. Khoảng
cách lớn ở tỉ lệ thực thi cho thấy lớp thứ hai còn yếu. Đây là điểm quan trọng
khi đánh giá LLM: một mô hình có thể đạt điểm cao trên các benchmark sinh mã
thuật toán, nơi tiêu chí là “qua bộ kiểm thử”, nhưng vẫn thất bại trong HPC, nơi
tiêu chí là “chạy đúng trên toolchain cụ thể với số tài nguyên cho trước”.

\subsection{Vì sao coroutine là điểm yếu}

Coroutine là biến thể có tỉ lệ lỗi cao nhất và khoảng cách lớn nhất so với mã con
người. Có ba lý do kỹ thuật.

Thứ nhất, coroutine yêu cầu người lập trình tự xây dựng phần hạ tầng: kiểu
\code{task}, \code{promise\_type}, scheduler, cách chờ và cách hủy. Không có một
“task chuẩn” trong C++20. Mỗi chương trình là một thiết kế riêng, nên mỗi lỗi
cũng là một lỗi thiết kế riêng, khó được che chắn bởi các mẫu quen thuộc.

Thứ hai, vòng đời frame là vấn đề ngữ nghĩa, không phải cú pháp. Việc
\code{final\_suspend} trả về \code{suspend\_never} hay \code{suspend\_always}
quyết định ai sở hữu frame, ai hủy nó, và \code{done()} có quan sát được trạng
thái kết thúc hay không. Trình biên dịch không thể cảnh báo vì cả hai lựa chọn
đều hợp lệ về mặt ngôn ngữ; chỉ ngữ cảnh sử dụng mới quyết định đúng sai.

Thứ ba, lỗi coroutine thường biểu hiện muộn: treo, segfault, hoặc ảnh sai. Chúng
không xuất hiện ở dòng khai báo mà ở vòng lặp chờ, ở hàm hủy, hoặc ở bước ghi ảnh.
Việc gỡ lỗi đòi hỏi hiểu luồng thực thi và quyền sở hữu tài nguyên — đúng những
thứ khó suy ra từ một prompt ngắn.

\subsection{“Có tham số” không đồng nghĩa “tôn trọng tham số”}

Một trong những bài học thực tiễn rõ nhất là không nên tin vào giao diện bề mặt.
Nhiều chương trình in ra hướng dẫn dòng lệnh đầy đủ, phân tích \code{-t}, nhưng
không dùng giá trị đó để giới hạn số luồng. Nguyên nhân có thể là:

\begin{itemize}
  \item Backend không phải OpenMP mà là TBB, nên biến môi trường bị bỏ qua.
  \item Số luồng được đặt nhưng thư viện tự quyết định dùng toàn bộ lõi.
  \item Tham số được truyền vào nhưng bị ghi đè bởi giá trị mặc định.
  \item Chương trình dùng song song lồng nhau, và cấp ngoài không bị giới hạn.
\end{itemize}

Cách duy nhất để xác nhận hành vi thật là \emph{đo}. Một phép đo mở rộng đơn giản
— chạy với 1, 2, 4, 8, 16, 20 lõi và vẽ thông lượng — sẽ lộ ra ngay chương trình
nào thực sự tôn trọng tham số. Đây là lý do bài báo không dừng ở kiểm tra biên
dịch mà còn đo hiệu năng.

\subsection{Chất lượng không tỉ lệ với số dòng mã}

Kết quả LOC và COCOMO cho thấy không có quan hệ đơn điệu giữa độ dài mã và chất
lượng. ChatGPT-5 sinh mã dài nhất ở bản phân tán nhưng cũng sạch nhất; LLaMA sinh
mã ngắn nhất nhưng phải thêm dòng và vẫn gặp lỗi thực thi. Mã ngắn có thể đạt được
bằng cách lược bỏ kiểm tra biên, bỏ xử lý lỗi hoặc giả định ngầm về dữ liệu đầu
vào — những thứ chỉ lộ ra khi chạy ở quy mô lớn.

COCOMO cũng cần được đọc thận trọng. Mô hình này ước lượng nỗ lực từ số dòng mã
và các hệ số kinh nghiệm, không mô hình hóa chi phí đồng bộ, truyền thông hay gỡ
lỗi song song. Vì vậy, một chương trình ít dòng hơn không chắc rẻ hơn về tổng chi
phí phát triển; nó chỉ có thể ít dòng hơn. Trong HPC, chi phí gỡ một lỗi
segfault do tranh chấp bộ nhớ có thể lớn hơn nhiều so với chi phí viết thêm vài
chục dòng kiểm tra.

\subsection{Hộp đen và quy mô mô hình}

Bài báo thẳng thắn thừa nhận khó khăn khi giải thích vì sao mô hình này tốt hơn
mô hình khác. LLM là hộp đen: dữ liệu huấn luyện, kiến trúc và chiến lược tối ưu
không được công bố đầy đủ. Có thể giả định dữ liệu huấn luyện chứa nhiều mã nguồn
công khai, nhưng không rõ có bao nhiêu mã song song chất lượng cao, bao nhiêu
thảo luận về TBB, coroutine hay MPI. Quy mô mô hình cũng là yếu tố mờ: các ước
lượng bên ngoài đặt ChatGPT-5 khoảng 600 tỉ tham số, Claude khoảng 175 tỉ, LLaMA
khoảng 400 tỉ, và một số mô hình nhỏ hơn khoảng 31 tỉ. Những con số này không đủ
để giải thích kết quả, vì chất lượng dữ liệu và căn chỉnh quan trọng không kém quy
mô.

Hệ quả phương pháp luận là không nên suy ra “mô hình lớn hơn luôn tốt hơn cho
HPC”. Cần đánh giá trên đúng bài toán, đúng toolchain và đúng tiêu chí. Một mô
hình nhỏ được tinh chỉnh trên mã song song có thể vượt mô hình lớn tổng quát
trong miền hẹp — đây cũng là động lực của các công trình như HPC-Coder, OMPGPT
hay MPIRigen.

\subsection{Hàm ý thực tiễn}

\subsubsection{Khi nào nên dùng LLM cho mã HPC}

Kết quả gợi ý ba mức sử dụng:

\begin{itemize}
  \item \textbf{Phù hợp}: tạo bản nháp cho các mẫu song song quen thuộc (vòng lặp
        OpenMP, async, thuật toán song song), đặc biệt khi bài toán đơn giản và
        có kiểm tra đầu ra rõ ràng.
  \item \textbf{Cần giám sát chặt}: mã coroutine, mã MPI có truyền thông phức
        tạp, mã có quản lý bộ nhớ thủ công. Đây là nơi lỗi ngữ nghĩa tập trung.
  \item \textbf{Chưa nên tin}: các tuyên bố về hiệu năng hoặc khả năng mở rộng
        nếu chưa đo. Không có cách nào thay thế phép đo thực nghiệm.
\end{itemize}

\subsubsection{Quy trình kiểm duyệt đề xuất}

Dựa trên các lỗi quan sát được, có thể áp dụng quy trình năm bước:

\begin{enumerate}
  \item \textbf{Biên dịch với cảnh báo}: bật \code{-Wall -Wextra}, coi cảnh báo
        là lỗi tiềm năng. Lỗi cú pháp là tầng dễ nhất nhưng vẫn cần loại bỏ.
  \item \textbf{Chạy trên dữ liệu nhỏ}: dùng ảnh nhỏ, ít vòng lặp để kiểm tra
        nhanh tính đúng trước khi đo hiệu năng.
  \item \textbf{Kiểm tra đầu ra}: so sánh ảnh với tham chiếu; kiểm tra định dạng
        tệp (P4/P5), kích thước, header.
  \item \textbf{Đo mở rộng}: chạy với nhiều mức lõi/nút, vẽ thông lượng; xác nhận
        tham số số luồng có hiệu lực thật.
  \item \textbf{Soát các điểm ngữ nghĩa}: vòng đời coroutine, quyền sở hữu bộ
        nhớ, kích thước buffer MPI, mốc đo thời gian.
\end{enumerate}

Bảng~\ref{tab:checklist} chuyển các bước này thành danh mục kiểm tra cụ thể.

\begin{table}[H]
\centering
\small
\caption{Danh mục kiểm tra mã song song do LLM sinh.}
\label{tab:checklist}
\begin{tabularx}{\textwidth}{@{}p{2.8cm} X@{}}
\toprule
\textbf{Nhóm} & \textbf{Điểm cần kiểm tra} \\
\midrule
Cú pháp & Biên dịch sạch với \code{-Wall -Wextra}; không dùng API không tồn tại; kiểu iterator đúng \\
Tham số & \code{-t}/\code{-p} được phân tích và \emph{sử dụng}; kiểm tra giá trị biên (0, 1, lớn hơn số lõi) \\
Số luồng & Xác nhận backend (TBB/OpenMP); dùng đúng API giới hạn; đo để chứng minh \\
Coroutine & \code{final\_suspend} phù hợp; \code{promise\_type} có \code{return\_void}/\code{return\_value}; không giải phóng hai lần \\
Bộ nhớ & Cấp phát/giải phóng cân xứng; kiểm tra biên; tránh chia sẻ giả \\
MPI & \code{sendcount}/\code{recvcount} khớp; kiểu dữ liệu khớp; root nhận đủ dữ liệu; xử lý trường hợp $P=1$ \\
Định dạng & Header và dữ liệu cùng định dạng (P4 với 1 bit, P5 với 8 bit); kích thước tệp hợp lệ \\
Đo đạc & Timer bao đúng phần tính toán; loại I/O; ghi rõ số lần lặp và cấu hình \\
Hiệu năng & Đường cong mở rộng có dốc; so sánh với mốc tuần tự; kiểm tra hiệu suất \\
\bottomrule
\end{tabularx}
\end{table}

\subsection{Hạn chế của nghiên cứu}

Có thể phân loại các hạn chế theo bốn mối đe dọa tính hợp lệ.

\subsubsection{Tính hợp lệ về cấu trúc}

Bài toán Mandelbrot là song song hiển nhiên, không có trao đổi dữ liệu giữa các
đơn vị tính toán. Vì vậy, nó không kiểm tra được các lỗi đồng bộ phức tạp, trao
đổi halo, collective nhiều bước hay phụ thuộc dữ liệu — những nơi LLM có thể sai
nhiều hơn. Kết quả về “mã MPI chạy được” chủ yếu nói về khung phân vùng và gom
dữ liệu, không nói về chất lượng của một bộ giải thực thụ.

\subsubsection{Tính hợp lệ nội tại}

Mỗi biến thể chỉ được sinh một lần với một prompt, không lặp lại nhiều lần, không
thay đổi nhiệt độ hay cách diễn đạt. Do đó không có phân tích thống kê về độ biến
thiên. Việc sửa lỗi bởi con người cũng tạo ra yếu tố chủ quan: mức độ “tối thiểu”
của một bản sửa có thể khác nhau giữa các lỗi. Ngoài ra, tiêu chí “ảnh đúng” được
kiểm tra trực quan, không có so sánh điểm ảnh tự động với ảnh tham chiếu.

\subsubsection{Tính hợp lệ bên ngoài}

Thực nghiệm chỉ dùng một trình biên dịch (GCC 11.5.0), một backend (TBB 2022.2),
một nền tảng CPU và một cụm. Kết quả có thể khác trên Clang, trên OpenMP thật,
trên GPU hoặc trên kiến trúc many-core. Bốn mô hình được đánh giá cũng chỉ là một
lát cắt tại thời điểm nghiên cứu; các mô hình mới như Gemini, DeepSeek-Coder hay
Grok chưa được xét.

\subsubsection{Tính hợp lệ của kết luận}

Một số kết luận dựa trên quan sát định tính về hình dạng đường cong, không có
kiểm định thống kê. Chính bài báo gốc tồn tại một số điểm không nhất quán giữa
bảng tổng hợp, bảng chi tiết, đồ thị và văn bản (đã tổng hợp đủ sáu điểm ở
Phần~\ref{sec:ketqua} và đối chiếu với Bảng~\ref{tab:detail}), nên việc trích dẫn
con số cụ thể cần đối chiếu đồng thời hai nguồn. Cuối cùng, COCOMO không được
thiết kế cho HPC, nên các suy luận về chi phí chỉ mang tính tương đối.
